\chapter{검증 결과(3): 적용 범위와 창 조정 모듈}

\section{팩 등급별 가치 지도(E18, E42)}

H5 기각은 사양 팩(640–920 V, 약 219s)에 대한 결론이다. 800 V급 주력 팩은 V\textsubscript{bat}/2가 400 V급 부품 범위 안에 있지만 900–1000 V급 팩은 그렇지 않으므로, 6장과 같은 규칙으로 셀 직렬수 180–240s × 부품 등급 B\textsubscript{420}·B\textsubscript{470} 전체를 다시 계산했다. 성립 조건은 계산 전에 고정했다. 정적 운전점의 버스가 [260 V, 부품 상한] 안에 있고, 조정 구간 정상점이 Stage B 포트 한계 안이며, 소자 직류 전압 사용률이 80 \% 이하인 것이다. 처음 계산(E18, 150 kHz)을 A25R과 최종 Stage C로 다시 계산했다(E36, E38).

\textbf{공정 기준으로 다시 계산했다(E42).} 이전 가치 지도는 ①의 보조 강압기(25 kW, 2위상)를 경부하 셰딩 없이 두었고 ④·R2-W의 Stage A는 설계된 셰딩을 썼으므로(②는 1위상이라 셰딩 대상 없음), 15.1절의 공정 기준(E41)과 같은 이상적 셰딩(동작점마다 1·2위상 중 손실이 작은 쪽)을 ①의 보조 강압기에 주고 다시 계산했다. 칸 판정 규칙은 실행 전에 정했다. 공정 기준 최소 손실 구조가 민감도 7종(자성체 × 0.5·× 2, 스위칭 × 0.7·× 1.5, 제어 손실 0, GaN 전이 × 2, ① 보조 강압기 25 kHz) 모두에서 최소이고 다음 구조보다 3 \% 넘게 작으면 “최선”, 아니면 “같은 수준”으로 적고 최선으로 주장하지 않는다. “같은 수준”의 후보는 각 경우에 (후보 손실 − 해당 경우 최소 손실)/후보 손실 ≤ 3 \%인 후보의 합집합이다. 표의 값은 E44 정정 뒤 값이다. ①의 보조 강압기 소자(C3M0016120K) E\textsubscript{oss} 127 µJ의 시험 전압을 데이터시트 조건(1000 V)으로 정정한 재계산(E44, 15.2절)에서 칸 분류는 같았고 B\textsubscript{470} 200·216s·사양 팩의 명목 최소가 ④에서 ①로 바뀌었다. 이때 판정 함수는 B\textsubscript{470} 200s의 같은 수준 후보에 R2-W410을 더했는데, 이는 함수가 각 경우의 실제 최소 구조가 아니라 명목 최소 구조(①)와 비교했기 때문이었다(자성체 × 2에서 실제 최소는 ④ 115.12 W, R2-W410 120.56 W, ① 128.80 W). 위 정의대로 각 경우의 실제 최소와 비교한 후속 재분류(E45, 판정 규칙을 실행 전에 기록)에서 R2-W410은 어느 경우에도 최소가 아니고 해당 경우 최소와의 차가 최소 4.51 \%여서 빠졌으며, 나머지 13칸의 분류와 후보는 같았다. E42–E44의 등록 판정과 원시 손실은 그대로 보존한다. 스위칭 배율을 Stage A 출력용량 하드 스위칭 항에도 준 재계산(E43)에서도 칸 판정은 같고 여유만 바뀌었다.

{\footnotesize\setlength{\tabcolsep}{3.5pt}
\begin{longtable}{L{\dimexpr 0.165\linewidth-2\tabcolsep\relax}L{\dimexpr 0.095\linewidth-2\tabcolsep\relax}L{\dimexpr 0.075\linewidth-2\tabcolsep\relax}L{\dimexpr 0.08\linewidth-2\tabcolsep\relax}L{\dimexpr 0.08\linewidth-2\tabcolsep\relax}L{\dimexpr 0.08\linewidth-2\tabcolsep\relax}L{\dimexpr 0.085\linewidth-2\tabcolsep\relax}L{\dimexpr 0.17\linewidth-2\tabcolsep\relax}L{\dimexpr 0.17\linewidth-2\tabcolsep\relax}}
\caption{팩 등급별 보조 경로 평균 손실과 칸 판정(최종 설계, 공정 기준)}\label{t:e18}\\
\toprule
\tabhead \textbf{팩(V\textsubscript{bat})} & \textbf{①} & \textbf{②} & \textbf{③} & \textbf{R2-W410} & \textbf{R2-W460} & \textbf{④(R2\newline 기본형)} & \textbf{B\textsubscript{420} 판정} & \textbf{B\textsubscript{470} 판정} \\
\midrule
\endfirsthead
\toprule
\tabhead \textbf{팩(V\textsubscript{bat})} & \textbf{①} & \textbf{②} & \textbf{③} & \textbf{R2-W410} & \textbf{R2-W460} & \textbf{④(R2\newline 기본형)} & \textbf{B\textsubscript{420} 판정} & \textbf{B\textsubscript{470} 판정} \\
\midrule
\endhead
\bottomrule
\endlastfoot
180s\newline(540–756 V) & 92.9\newline(104.1) & 134.0 & (불가) & = ④ & = ④ & 95.6 & ①·④ \mbox{같은 수준}\textsuperscript{a} & ①·④ \mbox{같은 수준}\textsuperscript{a} \\
192s\newline(576–806 V) & 95.1\newline(106.6) & 132.4 & (불가) & = ④ & = ④ & 96.7 & ①·④ \mbox{같은 수준} & ①·④ \mbox{같은 수준} \\
200s\newline(600–840 V) & 96.6\newline(108.3) & 133.5 & (불가) & 102.3 & = ④ & 97.6\newline(B\textsubscript{470}) & ①·R2-W410 \mbox{같은 수준} & ①·④ \mbox{같은 수준}\textsuperscript{b} \\
216s\newline(648–907 V) & 99.6\newline(111.8) & 134.0 & 144.4 & 126.5 & = ④ & 100.8\newline(B\textsubscript{470}) & \textbf{① 최선} & ①·④ \mbox{같은 수준} \\
사양\newline(640–920 V) & 100.2\newline(112.4) & 134.6 & 144.5 & 130.7 & = ④ & 101.5\newline(B\textsubscript{470}) & \textbf{① 최선} & ①·④ \mbox{같은 수준} \\
224s\newline(672–941 V) & 101.1\newline(113.5) & 135.6 & 144.5 & 139.5 & 104.0 & (불가) & \textbf{① 최선} & ①·R2-W460 \mbox{같은 수준} \\
240s\newline(720–1008 V) & (불가) & 137.8 & (불가) & (불가) & 129.6 & (불가) & ② 유일 성립 & R2-W460·② \mbox{같은 수준} \\
\end{longtable}
}
\srcnote{주: 단위 W, 주행 대표 프로파일의 대표 운전점 가중 평균(정적 모드 선택: 창 조정은 V\textsubscript{W} + 0.5 V 단일 경계, 셰딩은 위상 전류 기준), 공정 기준. ①의 괄호는 셰딩 없는 이전 값(E38)이다. “= ④”는 V\textsubscript{A}가 V\textsubscript{W}를 넘지 않아 Stage B를 두지 않는 경우, “④(B\textsubscript{470})”은 B\textsubscript{470}에서만 성립한다. 판정 칸에는 공정 기준에서 작은 구조를 먼저 적었다. 사양 팩의 R2-W410(130.7 W)은 가치 지도의 창 조정 계산 값이고, 6장 표의 ③′(130.4 W)은 같은 구조를 E15 방식으로 계산한 값이다. ① 열은 소자 조건 정정 뒤 값이다(E44). a: 180s ④는 400 V 부하의 정적 창 기준으로 성립한다. 최저 버스 268.77 V에서 48 V 서비스는 추가 확인이 필요하다(조정형 Stage C 설계 입력 하한 270 V. 192s의 최저 버스는 286.85 V). b: E44 판정 함수는 ①·④·R2-W410이었으나 경우별 실제 최소와 비교한 재분류(E45)에서 R2-W410이 빠졌다. 창 경계를 전압 상승 경로로 적용하면 R2-W410(99.1 W)도 같은 수준이 되고 하강 경로(105.5 W)에서는 표와 같다(운전 이력 의존, E46). 자료: 검증 근거 E18, E36, E38, E42, E43, E44, E45, E46.}

\begin{figure}[!htb]
\centering
\includegraphics[width=\linewidth]{\repo/docs/fig/report/fig7_1_value_map_2x2.png}
\caption{팩 등급 × 부품 등급의 구조 성립성과 손실(최종 설계, 공정 기준)}\label{f:e18}
\srcnote{주: (a) 비조정 부하 버스 전압 범위(경로·부하 강하 포함. 180s·셀 3.0 V의 반전압은 270.0 V, 그림의 하한은 약 268.77 V), (b)·(c) 부품 등급별 주행 프로파일 평균 보조 경로 손실(채운 표시 = 성립), (d) 500 V급 충전기 호환(셀 3.7 V·375 A의 직결 단자 V\textsubscript{A} 출력, R\textsubscript{A} 2.25 mΩ; E46에서 (d)만 다시 그림, 이전 5.55 mΩ 판보다 0.3 \% 작고 최대 237s는 같음). 자료: 검증 근거 E42, E44(소자 조건 정정 뒤), E46((d) 직결 단자 정의), E47((a) 축 이름).}
\end{figure}

결과(그림 \ref{f:e18})에서 다섯 가지를 확인했다.

\begin{itemize}
\item \textbf{③(전 범위 400 V 조정)은 216–224s와 사양 팩에서만 성립하고, 어느 팩에서도 최소 손실이 아니다.} 직렬 ±90 V로는 V\textsubscript{A} 약 310–490 V만 400 V로 맞출 수 있다.
\item \textbf{창 조정 모듈이 보조 손실에서 “최선”인 칸은 없다.} 공정 기준에서는 200s에서 ①이 R2-W410보다 5.6 \% 작고(민감도 −6.8 … +23.3 \%), 224s에서 ①이 R2-W460보다 2.8 \% 작다(−9.6 … +22.4 \%). 240s에서는 ①이 성립하지 않고 R2-W460이 ②보다 6.0 \% 작지만, 자성체 계수 × 0.5에서 0.4 \% 커진다. 세 칸 모두 “같은 수준”이다. 240s에서 R2-W460은 ②보다 스위치 VA가 22 kVA 크고 위치가 12개 많다.
\item \textbf{④(R2 기본형)와 ①은 800 V급 팩 전체에서 같은 수준이다.} 이전에 ④가 최선이던 B\textsubscript{470} 180–216s·사양 팩 칸은 공정 기준에서 ①이 명목상 1.1–2.9 \% 작지만 3 \% 안이고 민감도로 부호가 바뀐다(소자 조건 정정 전에는 차가 0.1–1.4 \%로 칸마다 순서가 달랐다). 15장의 ④ 대 ⑤ 결론과 같다.
\item \textbf{B\textsubscript{420}의 216·224s와 사양 팩에서는 ①이 모든 경우에 13.4 \% 이상 작다.} V\textsubscript{A}가 420 V를 넘는 팩에서 Stage A 기반 구조는 R2-W410(창 조정)이나 ②·③을 써야 하는데, 이들은 Stage A와 Stage B·조정기의 손실을 함께 진다.
\item \textbf{Stage A는 500 V급 충전기와도 호환된다.} 2:1 비는 1000 V급 팩을 500 V급 충전기 범위로 옮긴다. 충전기 직결 단자는 V\textsubscript{A}다(V\textsubscript{A} = V\textsubscript{bat}/2 + R\textsubscript{A}(I\textsubscript{ch} − i\textsubscript{path}), R\textsubscript{A} = 9.0 mΩ/4모듈 = 2.25 mΩ). 이 등가에서 375 A로 셀 4.2 V에 도달하는 최대 직렬수는 237s다(240s는 셀 약 4.16 V까지; 5.25·5.55 mΩ 정의로도 같음). 무보조 375 A에서 출력 상한 400·500 V 충전기가 도달하는 팩 전압은 798.31·998.31 V(500 V·240s의 375 A 상한 4.1596 V/셀)이며, 정전류·정전압 종료, BMS 전압 기준, 충전 시간, 차량 부하에 따른 충전 가능 구간은 계산에 없다. 셀 3.7 V·375 A·보조부하 0의 단자 출력은 125.19–166.82 kW(180–240s)다. 216s는 직결 단자 기준 375 A에서 150.166 kW, 150 kW에 374.585 A(8위상)이고, (V\textsubscript{bat}/2)I 기준으로는 375 A에서 149.850 kW라 150 kW에 375.375 A(10위상)가 필요하다(E46; 이전 판의 150.588 kW·373.542 A는 V\textsubscript{A} 앞에 별도 3 mΩ 경로를 둔 계산). 팩별 150 kW 충전 경로 손실은 (V\textsubscript{bat}/2)I 기준으로 위상 수를 골라(위상당 ≤ 46.875 A, 2위상 단위) 계산했다: 180–216s는 10위상(5모듈, 375–450 A)으로 590–632 W, 224s·240s·사양 팩은 8위상(338–370 A)으로 565–582 W다. Stage A 공진 인덕터의 저장에너지(8위상)는 18.9 mJ로 150 kW 승압기(약 3.3 J)의 약 1/175이다.
\end{itemize}

\textbf{Q3(적용 범위)의 답:} 전 범위 조정으로는 이점이 없다. 창 조정 모듈은 V\textsubscript{A}가 부품 상한을 넘는 팩(B\textsubscript{420} 200s, B\textsubscript{470} 224·240s)에서 400 V급 부품을 쓰게 하는 \textbf{성립 수단}이며, 보조 손실은 그 칸의 다른 성립 구조와 같은 수준이므로 이점으로 주장하지 않는다. ① 대비로 남는 차이는 충전 경로 손실(사양 팩 375 A 충전에서 전용 승압기(16–50 kHz)의 0.32–0.51배)과 자성체 저장에너지이고, ② 대비로는 지표상 이점이 없다. 기여는 손실이 아니라 성립 조건과 고장 강등 순서(9장)다. 산업 방식 대비 위치는 15장에서 다룬다.

\section{창 조정 모듈 검증(E19)}

Stage B가 필요한 대표 두 사례(R2-W460: 240s, 720–1008 V, B\textsubscript{470} 부품. R2-W410: 사양 팩, 640–920 V, B\textsubscript{420} 부품)를 평균모델과 채택 제어기로 시험했다. 판정은 E13의 Tier 기준에 부품 등급 상한을 적용했다. 이 시험과 회로 시험(TB08)은 Stage A를 150 kHz 설계의 출력저항(경로 저항 13.2 mΩ에 포함)을 가진 전압원으로 두었고, A25R로 다시 시험하지 않았다. A25R의 경로 저항(12.0 mΩ)은 시험의 공차 범위(±10 \%) 안이지만 V\textsubscript{A} 노드 커패시터(2.4 mF)는 모델에 없다.

처음 정한 DAB 정격(정적 최대 링크전력의 1.25배)은 최대 배터리 전압의 Tier-2(1→25 kW) 스텝에서 포화되어 순시 전력이 정격을 넘었다(1.25배 정격 R2-W460 3.0 kW 71/72, R2-W410 3.5 kW 70/72. 이전 판의 68/72는 R2-W410 3.0 kW 행이었다). 포화 없는 과도 첨두가 정적 최대의 1.77–1.78배여서 정격 규칙을 1.8배로 바꿨다(표 \ref{t:e19}).

{\footnotesize
\begin{longtable}{L{\dimexpr 0.2982\linewidth-2\tabcolsep\relax}L{\dimexpr 0.3509\linewidth-2\tabcolsep\relax}L{\dimexpr 0.3509\linewidth-2\tabcolsep\relax}}
\caption{창 조정 모듈의 검증 결과(평균모델, 최종 정격)}\label{t:e19}\\
\toprule
\tabhead \textbf{시험} & \textbf{R2-W460(4.0 kW)} & \textbf{R2-W410(5.0 kW)} \\
\midrule
\endfirsthead
\toprule
\tabhead \textbf{시험} & \textbf{R2-W460(4.0 kW)} & \textbf{R2-W410(5.0 kW)} \\
\midrule
\endhead
\bottomrule
\endlastfoot
Tier-1·Tier-2(공차 16종, 지연 +3.0 µs 포함) & 456/456, 최대 81.65 A, 버스 455.7–464.3 V & 264/264, 최대 85.35 A, 버스 405.1–414.8 V \\
켬(영벡터 → 창 조정, 램프 0–2 ms × 부하 3종) & 12/12, 최대 53.9 A, 버스 ≤ 465.0 V & 12/12, 최대 60.4 A, 버스 ≤ 415.0 V \\
고 SOC 기동 인계 & 9/9, 최대 77.9 A & 9/9, 최대 84.8 A \\
끔(→ 바이패스) & 6/6, 바이패스 직전 직렬 ≤ 0.17 V & 6/6, ≤ 0.19 V \\
DAB 고장 → 기준 램프 1 ms → 바이패스 & 2/2, 경로 ≤ 69.8 A, 링크 ≤ 109.7 V & 2/2, 경로 ≤ 78.2 A, 링크 ≤ 112.2 V \\
같은 대응, 램프 0.5 ms & 1/2(25 kW에서 125.1 A) & 1/2(220.2 A) \\
H-브리지 전 소자 차단 → 50 µs 후 바이패스 & 0/2, 최대 740.5 A & 0/2, 838.2 A \\
\end{longtable}
}

\textbf{DAB 고장은 Stage A를 멈추지 않고 강등할 수 있다.} 링크가 직렬 포트의 전력을 흡수하는 동안 직렬전압을 0으로 줄인 뒤 바이패스한다(③에서는 배터리 840 V 초과에서 Stage A 협조 정지가 필요했다). 강등 뒤 버스는 V\textsubscript{A}(503.9·459.9 V)로 부품 상한을 넘으므로 400 V 부하는 분리해야 하며, 조정형 Stage C가 48 V 공급을 이어 간다. \textbf{H-브리지 전 소자 차단 고장에서는 바로 바이패스하면 안 된다.} 다이오드가 직렬 커패시터를 링크 쪽으로 충전한 상태(직렬 45.8–52.1 V)에서 단락하기 때문이며, 강등 순서는 9장에서 정했다.

\section{창 조정 모듈의 회로 확인(TB08)}

R2-W460을 소자 수준 회로에서 다시 확인했다. 기존 Stage B 회로(TB04B)에서 DAB 권선비(4.6:1)와 누설 인덕턴스(14.68 µH)만 바꾸고 바이패스 스위치, 모드 시퀀서, 제어기 무충격 재초기화를 더했으며, 다섯 사례 모두 통과했다(TB08, 5/5). 창 조정 Tier-2(1008 V, 1→25 kW)는 버스 455.95–460.79 V, 경로 69.66 A(스위칭 ODE 대비 +1.84 A, 허용 2.03 A), 링크전력 3.22 kW, 격자 0.625 ns 리플 0.050 V였다. Tier-1(930 V, 15→25 kW)은 버스 458.31–460.73 V(DAB 저전류 영역에서도 링크 99.94–100.02 V), 끔(910 V, 25 kW)은 경로 ≤ 57.2 A·바이패스 직전 직렬 −0.22 V(평균모델 −0.17 V), DAB 고장 강등(1008 V, 25 kW)은 링크 최대 109.76 V(평균모델 109.72 V)·경로 ≤ 69.7 A·강등 뒤 버스 503.2 V, 켬(930 V, 25 kW)은 경로 ≤ 73.4 A·버스 ≤ 465.4 V·최종 460.02 V였다.


회로에서 새로 드러난 두 현상은 모두 한계 안이다(그림 \ref{f:tb08}). 영벡터에서도 직렬 포트에 0.2–0.6 V가 남아 바이패스를 닫는 순간 직렬 커패시터가 방전되고(최대 −173 A, 약 1 µs, 소자당 펄스 정격 329 A\cite{b15}의 0.26배), 바이패스를 열면 경로전류 40 \%가 영벡터 경로로 옮겨 가며 19.6 kHz로 출렁인다(경로 ≤ 73.4 A). 둘째 현상은 보호 상태기계의 켬 순서(전류 조향)에서 없앴다(9.1절).

\begin{figure}[!htb]
\centering
\includegraphics[width=15.5cm]{\repo/docs/fig/tb08_bypass_transitions.png}
\caption{TB08 바이패스 투입·열림 순시 파형}\label{f:tb08}
\srcnote{주: (a) DAB 고장 강등의 바이패스 투입, (b) 켬 전환의 바이패스 열림. 기생 인덕턴스가 없는 이상 루프 값이다.}
\end{figure}

\section{측정 사양을 넣은 창 조정 확인(E43)}

창 조정 모듈의 검증(E19, TB08)은 이상 측정으로 했는데, 처음 구조의 버스 센서 창(350–450 V, E12)으로는 R2-W460의 버스(460 V, 강등 뒤 최대 504 V)를 측정할 수 없다(외부 검토 2026-11-05). 그래서 측정 사양 후보 S1(버스 0–600 V 단일 채널, 값은 표 \ref{t:req})에 범위 밖 포화를 넣은 센서로 스위칭 모델에서 다시 판정했다(기준과 채택 규칙은 실행 전에 정함). R2-W460·R2-W410 각 3개 배터리 전압, 공칭·지연 +3.0 µs, Tier-1 10계단·Tier-2 2계단의 144건을 이상 측정, S1 잡음·양자화, S1 + 오차(두 방향)의 네 변형으로 계산하고, 하드웨어 버스 과전압 비교기(문턱 = 상한)와 버스 조정 상실 판단(|측정 − V\textsubscript{W}| > 20 V가 0.5 ms)의 오동작도 확인했다.


결과는 네 변형 모두 144/144 통과였다. S1 변형의 마지막 1 ms 리플(PWM 주기평균의 변동)은 최대 0.057 V(기준 0.1 V), 순시 경로·브리지 전류는 90.6 A(기준 98 A)였고, 이득·오프셋 오차로 실제 버스는 R2-W460에서 457.6–462.4 V로 정착했다(최대 466.7 V, R2-W410 417.0 V). 비교기 쪽 측정 최대는 464.5 V(R2-W410 414.9 V)로 잘못된 차단이 없었고 조정 상실 판단도 걸리지 않았다. 최종 설계 사건의 V\textsubscript{A}·버스 최대 516.3 V(W460 클램프, 문턱 +1 \%)의 1.1배(568 V)를 600 V 범위가 덮으므로 S1을 측정 사양으로 채택했다. 한계는 센서 대역폭을 따로 넣지 않은 점, 보호 오동작 확인이 측정 오차를 실제 궤적에 더한 보수 계산이라는 점이며, 창 조정 모듈의 고장 강등(E17·E19–E21)은 이 사양으로 다시 계산하지 않았다.
